% ══════════════════════════════════════════════════════════════════════════════
%                         فصل دوم
%                 فلسفه و معرفت‌شناسی چپ و راست
% ══════════════════════════════════════════════════════════════════════════════

\chapter{فلسفه و معرفت‌شناسی چپ و راست}
\label{ch:philosophy}

\begin{flushright}
\textit{«تفاوت چپ و راست، پیش از آنکه سیاسی باشد، فلسفی است.»}
\end{flushright}

% ══════════════════════════════════════════════════════════════════════════════
%                         مقدمه فصل
% ══════════════════════════════════════════════════════════════════════════════

\lettrine[lines=3, lraise=0.1, nindent=0.5em]{\textcolor{RightBlue}{آ}}{یا} 
چپ و راست صرفاً سلیقه‌های سیاسی متفاوت‌اند، یا ریشه در جهان‌بینی‌های عمیق‌تری دارند؟ در این فصل نشان خواهیم داد که پشت اختلافات ظاهراً سیاسی، تفاوت‌های بنیادین فلسفی نهفته است: تفاوت در نگاه به انسان، تاریخ، عقل، عدالت و آزادی. درک این مبانی فلسفی، کلید فهم عمیق‌تر سیاست است.

% ══════════════════════════════════════════════════════════════════════════════
\section{انسان‌شناسی فلسفی: انسان چیست؟}
% ══════════════════════════════════════════════════════════════════════════════

بنیادی‌ترین اختلاف فلسفی میان چپ و راست، در پاسخ به این پرسش نهفته است: \textbf{طبیعت انسان چیست؟}

\subsection{نگاه بدبینانه: انسان ناقص}

سنت محافظه‌کار و بخش‌هایی از راست، بر \keyword{نقصان ذاتی انسان} تأکید می‌کنند. این دیدگاه ریشه در الهیات مسیحی (گناه نخستین) و فلسفه سیاسی هابز دارد.

\begin{personbox}[توماس هابز (۱۵۸۸-۱۶۷۹)]

\textbf{ملیت:} انگلیسی

\textbf{اثر کلیدی:} \book{لویاتان} (۱۶۵۱)

\textbf{ایده محوری:} وضع طبیعی به مثابه جنگ همه علیه همه

\mediumspace

\person{هابز} استدلال کرد که انسان‌ها در «وضع طبیعی» (بدون دولت) در جنگ دائمی زندگی می‌کنند. زندگی در این وضعیت «تنها، فقیرانه، کثیف، حیوانی و کوتاه» است. از این رو، انسان‌ها برای صلح و امنیت، قدرت مطلق را به «لویاتان» (دولت) واگذار می‌کنند.

\mediumspace

\textbf{نقل‌قول معروف:} \textit{«انسان گرگ انسان است.»}

\mediumspace

\textbf{پیامد سیاسی:} نیاز به قدرت قاهره برای کنترل طبیعت خودخواه بشر
\end{personbox}

\begin{defbox}[تعریف: بدبینی انسان‌شناختی]
دیدگاهی که معتقد است:
\begin{itemize}[nosep]
    \item انسان ذاتاً خودخواه، پرخاشگر یا ناقص است
    \item این نقصان قابل «درمان» نیست، فقط قابل «مهار» است
    \item نهادها (دولت، مذهب، سنت) برای کنترل این طبیعت ضروری‌اند
    \item آرمان‌شهرها محکوم به شکست‌اند
\end{itemize}
\end{defbox}

\subsection{نگاه خوش‌بینانه: انسان قابل تکامل}

در مقابل، سنت روشنگری و بخش‌های مهمی از چپ، بر \keyword{قابلیت تکامل انسان} تأکید می‌کنند. شاخص‌ترین نماینده این دیدگاه، ژان ژاک روسو است.

\begin{personbox}[ژان ژاک روسو (۱۷۱۲-۱۷۷۸)]

\textbf{ملیت:} سوئیسی-فرانسوی

\textbf{آثار کلیدی:} \book{قرارداد اجتماعی} (۱۷۶۲)، \book{امیل} (۱۷۶۲)

\textbf{ایده محوری:} «انسان آزاد زاده می‌شود، اما همه‌جا در زنجیر است»

\mediumspace

\person{روسو} استدلال کرد که انسان در وضع طبیعی «وحشی نجیب» بود—صلح‌جو و خوشبخت. این جامعه و نهادهایش هستند که انسان را فاسد کرده‌اند. مالکیت خصوصی، نابرابری و تمدن، منشأ بدبختی بشرند.

\mediumspace

\textbf{نقل‌قول معروف:} \textit{«نخستین کسی که زمینی را محصور کرد و گفت "این مال من است"، بنیان‌گذار واقعی جامعه مدنی بود.»}

\mediumspace

\textbf{پیامد سیاسی:} با تغییر نهادها، می‌توان انسان بهتری ساخت
\end{personbox}

\begin{figure}[htbp]
\centering
\begin{tikzpicture}[scale=0.9]

% عنوان
\node[font=\large\bfseries, text=DarkGray] at (0, 6.5) {\fa{دو نگاه به طبیعت انسان}};

% کادر هابز (راست)
\fill[RightBlue, opacity=0.2, rounded corners=10pt] (-7, -0.5) rectangle (-1, 5.5);
\draw[line width=2pt, color=RightBlue, rounded corners=10pt] (-7, -0.5) rectangle (-1, 5.5);

\node[text=RightBlue, font=\bfseries\large] at (-4, 5) {\fa{هابز / محافظه‌کاری}};

\node[text=DarkGray, font=\small, text width=5cm, align=center] at (-4, 3.8) {
    \fa{انسان ذاتاً خودخواه}
};
\node[text=DarkGray, font=\small, text width=5cm, align=center] at (-4, 3) {
    \fa{وضع طبیعی = جنگ}
};
\node[text=DarkGray, font=\small, text width=5cm, align=center] at (-4, 2.2) {
    \fa{نیاز به کنترل بیرونی}
};
\node[text=DarkGray, font=\small, text width=5cm, align=center] at (-4, 1.4) {
    \fa{نهادها برای مهار}
};
\node[text=RightBlueDark, font=\small\bfseries, text width=5cm, align=center] at (-4, 0.3) {
    \fa{نتیجه: محدودیت ضروری است}
};

% کادر روسو (چپ)
\fill[LeftRed, opacity=0.2, rounded corners=10pt] (1, -0.5) rectangle (7, 5.5);
\draw[line width=2pt, color=LeftRed, rounded corners=10pt] (1, -0.5) rectangle (7, 5.5);

\node[text=LeftRed, font=\bfseries\large] at (4, 5) {\fa{روسو / رادیکالیسم}};

\node[text=DarkGray, font=\small, text width=5cm, align=center] at (4, 3.8) {
    \fa{انسان ذاتاً نیک}
};
\node[text=DarkGray, font=\small, text width=5cm, align=center] at (4, 3) {
    \fa{وضع طبیعی = صلح}
};
\node[text=DarkGray, font=\small, text width=5cm, align=center] at (4, 2.2) {
    \fa{جامعه فاسد می‌کند}
};
\node[text=DarkGray, font=\small, text width=5cm, align=center] at (4, 1.4) {
    \fa{نهادها برای رهایی}
};
\node[text=LeftRedDark, font=\small\bfseries, text width=5cm, align=center] at (4, 0.3) {
    \fa{نتیجه: تغییر ممکن است}
};

% علامت مقابله
\node[text=GoldAccent, font=\Huge\bfseries] at (0, 2.5) {\fa{⚔}};

\end{tikzpicture}
\caption{مقایسه انسان‌شناسی هابز و روسو}
\label{fig:hobbes-rousseau}
\end{figure}

\begin{debatebox}[هابز در برابر روسو: کدام درست است؟]
\textbf{طرفداران هابز:}
\begin{itemize}[nosep]
    \item شواهد تاریخی: جوامع بدون دولت اغلب خشونت‌بار بوده‌اند
    \item روان‌شناسی تکاملی: انسان برای بقا رقابت می‌کند
    \item تجربه قرن بیستم: آرمان‌شهرها به کابوس تبدیل شدند
\end{itemize}

\textbf{طرفداران روسو:}
\begin{itemize}[nosep]
    \item جوامع شکارچی-گردآورنده نسبتاً برابری‌طلب بودند
    \item خشونت سازمان‌یافته محصول تمدن است، نه طبیعت
    \item «طبیعت» انسان خود ساخته اجتماعی است
\end{itemize}

\textbf{حقیقت شاید جایی میان این دو باشد.}
\end{debatebox}

\subsection{جان لاک: راه میانه}

\begin{personbox}[جان لاک (۱۶۳۲-۱۷۰۴)]

\textbf{ملیت:} انگلیسی

\textbf{آثار کلیدی:} \book{دو رساله در باب حکومت} (۱۶۸۹)، \book{رساله در باب فهم بشری}

\textbf{ایده محوری:} حقوق طبیعی و حکومت محدود

\mediumspace

\person{لاک} نه به اندازه هابز بدبین بود و نه به اندازه روسو خوش‌بین. او معتقد بود انسان در وضع طبیعی از \textbf{حقوق طبیعی} (زندگی، آزادی، مالکیت) برخوردار است و عقل، راهنمای اوست. اما چون برخی این حقوق را نقض می‌کنند، نیاز به حکومت هست—حکومتی محدود و بر اساس رضایت.

\mediumspace

\textbf{پیامد:} پدر لیبرالیسم کلاسیک
\end{personbox}

% ══════════════════════════════════════════════════════════════════════════════
\section{معرفت‌شناسی سیاسی: چگونه می‌دانیم؟}
% ══════════════════════════════════════════════════════════════════════════════

اختلاف دیگر، در این است که \textbf{منبع شناخت سیاسی} چیست؟ از کجا بدانیم جامعه خوب چگونه است؟

\subsection{خرد جمعی در برابر خرد فردی}

\begin{comparebox}[دو منبع شناخت سیاسی]
\begin{table}[H]
\centering
\begin{tabular}{r|c|c}
\toprule
\textbf{موضوع} & \textbf{محافظه‌کاری} & \textbf{روشنگری/چپ} \\
\midrule
منبع شناخت & سنت و تجربه انباشته & عقل فردی و علم \\
روش & آزمون‌وخطای تاریخی & تحلیل عقلانی \\
نگاه به گذشته & مخزن حکمت & مانع پیشرفت \\
نگاه به تغییر & محتاطانه و تدریجی & جسورانه و بنیادین \\
اعتماد به نخبگان & متخصصان سنت & متخصصان علم \\
\bottomrule
\end{tabular}
\end{table}
\end{comparebox}

\begin{quotebox}[ادموند برک]
\textit{«ما از عقل فردی خود می‌ترسیم، زیرا گمان می‌کنیم سرمایه عقل در هر انسان اندک است. بهتر است از انبار عمومی عقل ملت‌ها و قرون بهره ببریم.»}
\end{quotebox}

\begin{defbox}[تعریف: عقل‌گرایی سیاسی]
باور به اینکه:
\begin{itemize}[nosep]
    \item جامعه را می‌توان بر اساس اصول عقلی (دوباره) طراحی کرد
    \item سنت صرفاً به این دلیل که قدیمی است، ارزشمند نیست
    \item پیشرفت از طریق به‌کارگیری خرد در سیاست ممکن است
\end{itemize}
این نگرش در روشنگری و چپ مدرن غالب است.
\end{defbox}

\subsection{نقد ایدئولوژی: مارکس و آگاهی کاذب}

\person{کارل مارکس} معرفت‌شناسی سیاسی را وارد مرحله جدیدی کرد. او پرسید: چرا مردم به ایده‌هایی باور دارند که به ضررشان است؟

\begin{defbox}[تعریف: ایدئولوژی و آگاهی کاذب]
در سنت مارکسی:
\begin{itemize}[nosep]
    \item \textbf{ایدئولوژی:} مجموعه‌ای از ایده‌ها که منافع طبقه حاکم را توجیه می‌کند
    \item \textbf{آگاهی کاذب:} وقتی طبقات فرودست، ایدئولوژی طبقه حاکم را درونی می‌کنند
    \item \textbf{نقد ایدئولوژی:} آشکارسازی منافع پنهان پشت ایده‌های «بی‌طرف»
\end{itemize}
\end{defbox}

\begin{quotebox}[کارل مارکس]
\textit{«ایده‌های حاکم در هر دوره، ایده‌های طبقه حاکم‌اند.»}
\end{quotebox}

این رویکرد، هم قوت‌هایی دارد (توجه به رابطه قدرت و دانش) و هم ضعف‌هایی (آیا نقد ایدئولوژی خود ایدئولوژیک نیست؟).

\subsection{مایکل اوکشات: نقد عقل‌گرایی از راست}

\begin{personbox}[مایکل اوکشات (۱۹۰۱-۱۹۹۰)]

\textbf{ملیت:} انگلیسی

\textbf{اثر کلیدی:} \book{عقل‌گرایی در سیاست} (۱۹۶۲)

\textbf{ایده محوری:} تمایز میان «دانش فنی» و «دانش عملی»

\mediumspace

\person{اوکشات} استدلال کرد که عقل‌گرایان سیاسی خطای بزرگی مرتکب می‌شوند: آنان گمان می‌کنند سیاست مانند مهندسی است و می‌توان جامعه را طبق نقشه ساخت. اما سیاست بیشتر شبیه آشپزی است—نیاز به «دانش عملی» دارد که از تجربه و سنت می‌آید، نه از کتاب.

\mediumspace

\textbf{نقل‌قول:} \textit{«در سیاست، مردان اغلب مانند مردی عمل می‌کنند که تنها با کتاب آشپزی آشپزی کند.»}
\end{personbox}

% ══════════════════════════════════════════════════════════════════════════════
\section{فلسفه تاریخ: تاریخ به کجا می‌رود؟}
% ══════════════════════════════════════════════════════════════════════════════

نگاه به تاریخ نیز چپ و راست را از هم جدا می‌کند.

\subsection{پیشرفت‌گرایی: روایت غالب چپ}

\begin{defbox}[تعریف: پیشرفت‌گرایی]
باور به اینکه:
\begin{itemize}[nosep]
    \item تاریخ جهت‌دار است و به سوی بهبود حرکت می‌کند
    \item عقل و علم، موتور پیشرفت‌اند
    \item آینده بهتر از گذشته خواهد بود
    \item موانع پیشرفت (سنت، خرافه، نابرابری) باید برداشته شوند
\end{itemize}
این ایده از روشنگری آمده و در چپ مدرن تأثیرگذار بوده است.
\end{defbox}

\begin{figure}[htbp]
\centering
\begin{tikzpicture}[scale=0.85]

% عنوان
\node[font=\large\bfseries, text=DarkGray] at (0, 7) {\fa{سه روایت از تاریخ}};

% ─────────── روایت ۱: پیشرفت‌گرایی (چپ) ───────────
\node[text=LeftRed, font=\bfseries] at (-5, 5.5) {\fa{پیشرفت‌گرایی}};

\draw[line width=2pt, color=LeftRed, -{Stealth[length=8pt]}] 
    (-7, 4) -- (-3, 5);
    
\node[text=MediumGray, font=\footnotesize] at (-7.3, 3.7) {\fa{گذشته}};
\node[text=MediumGray, font=\footnotesize] at (-2.5, 5.2) {\fa{آینده}};
\node[text=LeftRed, font=\small] at (-5, 3.5) {\fa{صعودی}};

% ─────────── روایت ۲: چرخه‌ای (برخی محافظه‌کاران) ───────────
\node[text=RightBlue, font=\bfseries] at (0, 5.5) {\fa{چرخه‌ای}};

\draw[line width=2pt, color=RightBlue] 
    (0, 4) circle (0.8);
\draw[line width=2pt, color=RightBlue, -{Stealth[length=6pt]}] 
    (0.7, 4.3) arc (30:-60:0.8);

\node[text=RightBlue, font=\small] at (0, 3) {\fa{تکرار}};

% ─────────── روایت ۳: زوال (بدبینانه) ───────────
\node[text=MediumGray, font=\bfseries] at (5, 5.5) {\fa{زوال‌گرایی}};

\draw[line width=2pt, color=MediumGray, -{Stealth[length=8pt]}] 
    (3, 5) -- (7, 4);

\node[text=MediumGray, font=\footnotesize] at (2.7, 5.2) {\fa{عصر طلایی}};
\node[text=MediumGray, font=\footnotesize] at (7.3, 3.7) {\fa{انحطاط}};
\node[text=MediumGray, font=\small] at (5, 3.5) {\fa{نزولی}};

% ─────────── روایت مارکس ───────────
\node[text=LeftRedDark, font=\bfseries] at (0, 1.5) {\fa{ماتریالیسم تاریخی مارکس}};

\draw[line width=2pt, color=LeftRedDark] (-6, 0) -- (-3.5, 0);
\draw[line width=2pt, color=LeftRedDark] (-3.5, 0) -- (-3.5, 0.5) -- (-1, 0.5);
\draw[line width=2pt, color=LeftRedDark] (-1, 0.5) -- (-1, 1) -- (1.5, 1);
\draw[line width=2pt, color=LeftRedDark] (1.5, 1) -- (1.5, 1.5) -- (4, 1.5);
\draw[line width=2pt, color=LeftRedDark, -{Stealth[length=8pt]}] (4, 1.5) -- (4, 2) -- (6, 2);

\node[text=LeftRedDark, font=\tiny] at (-5, -0.4) {\fa{فئودالیسم}};
\node[text=LeftRedDark, font=\tiny] at (-2.2, 0.1) {\fa{سرمایه‌داری}};
\node[text=LeftRedDark, font=\tiny] at (0.2, 0.6) {\fa{سوسیالیسم}};
\node[text=LeftRedDark, font=\tiny] at (2.7, 1.1) {\fa{کمونیسم}};
\node[text=LeftRedDark, font=\tiny] at (6.3, 2) {\fa{پایان}};

\node[text=LeftRedDark, font=\small] at (0, -1) {\fa{پلکانی با انقلاب‌ها}};

\end{tikzpicture}
\caption{روایت‌های مختلف از سیر تاریخ}
\label{fig:history-views}
\end{figure}

\subsection{تاریخ‌گرایی مارکسی}

مارکس نظریه‌ای علمی‌گونه از تاریخ ارائه داد: \term{ماتریالیسم تاریخی}.

\begin{historybox}[ماتریالیسم تاریخی در یک نگاه]
\begin{enumerate}[nosep]
    \item تاریخ، تاریخ مبارزه طبقاتی است
    \item هر دوره تاریخی، شیوه تولید خاصی دارد (برده‌داری، فئودالیسم، سرمایه‌داری)
    \item تضادهای درونی هر شیوه، آن را نابود و شیوه جدید را جایگزین می‌کند
    \item سرمایه‌داری نیز محکوم به فروپاشی است
    \item پرولتاریا (طبقه کارگر) آخرین طبقه انقلابی است
    \item نقطه پایان: جامعه بی‌طبقه (کمونیسم)
\end{enumerate}
\end{historybox}

\subsection{نقد محافظه‌کارانه: تاریخ معلم نیست، هشدار است}

محافظه‌کاران معمولاً نسبت به هر دو دیدگاه (پیشرفت قطعی یا فروپاشی قطعی) بدبین‌اند.

\begin{quotebox}[رابرت نیسبت، جامعه‌شناس محافظه‌کار]
\textit{«ایده پیشرفت یکی از خطرناک‌ترین ایده‌های تاریخ بشر است. با وعده آینده درخشان، حال را قربانی می‌کند.»}
\end{quotebox}

% ══════════════════════════════════════════════════════════════════════════════
\section{عدالت و برابری}
% ══════════════════════════════════════════════════════════════════════════════

مفهوم \keyword{عدالت} شاید مناقشه‌برانگیزترین مفهوم فلسفه سیاسی باشد.

\subsection{عدالت چیست؟}

\begin{defbox}[تعریف: انواع عدالت]
\begin{itemize}[nosep]
    \item \textbf{عدالت توزیعی:} چگونگی تقسیم کالاها و فرصت‌ها در جامعه
    \item \textbf{عدالت تصحیحی:} جبران خطاها و جرائم
    \item \textbf{عدالت رویه‌ای:} عادلانه بودن فرآیندها (دادگاه عادل، انتخابات عادل)
    \item \textbf{عدالت بین‌نسلی:} مسئولیت ما در برابر نسل‌های آینده
\end{itemize}
اختلاف چپ و راست عمدتاً درباره \textbf{عدالت توزیعی} است.
\end{defbox}

\subsection{جان رالز: عدالت به مثابه انصاف}

\begin{personbox}[جان رالز (۱۹۲۱-۲۰۰۲)]

\textbf{ملیت:} آمریکایی

\textbf{اثر کلیدی:} \book{نظریه‌ای در باب عدالت} (۱۹۷۱)

\textbf{ایده محوری:} پرده جهل و اصول عدالت

\mediumspace

\person{رالز} پرسید: اگر ما نمی‌دانستیم در جامعه چه جایگاهی خواهیم داشت (پشت «پرده جهل»)، چه اصولی را برای جامعه انتخاب می‌کردیم؟

\textbf{پاسخ رالز:}
\begin{enumerate}[nosep]
    \item \textbf{اصل آزادی:} بیشترین آزادی‌های پایه برای همه
    \item \textbf{اصل تفاوت:} نابرابری‌ها فقط وقتی موجه‌اند که به نفع محروم‌ترین‌ها باشند
\end{enumerate}

\mediumspace

\textbf{پیامد:} دفاع فلسفی از دولت رفاه لیبرال
\end{personbox}

\subsection{رابرت نوزیک: عدالت به مثابه استحقاق}

\begin{personbox}[رابرت نوزیک (۱۹۳۸-۲۰۰۲)]

\textbf{ملیت:} آمریکایی

\textbf{اثر کلیدی:} \book{آنارشی، دولت و آرمان‌شهر} (۱۹۷۴)

\textbf{ایده محوری:} حقوق مالکیت و دولت حداقلی

\mediumspace

\person{نوزیک} در پاسخ به رالز استدلال کرد که عدالت درباره \textbf{فرآیند} است، نه \textbf{الگوی توزیع}. اگر من چیزی را به‌طور عادلانه کسب کرده‌ام (کار، مبادله، هدیه)، آن متعلق به من است و دولت حق بازتوزیع آن را ندارد.

\textbf{استدلال معروف:} اگر همه با ویلت چمبرلین (بسکتبالیست) قرارداد ببندند که ۲۵ سنت از هر بلیط به او برسد، آیا ثروت حاصل ناعادلانه است؟

\mediumspace

\textbf{پیامد:} دفاع از دولت حداقلی و لیبرتارینیسم
\end{personbox}

\begin{figure}[htbp]
\centering
\begin{tikzpicture}[scale=0.85]

% عنوان
\node[font=\large\bfseries, text=DarkGray] at (0, 6) {\fa{رالز در برابر نوزیک}};

% محور
\draw[line width=3pt, MediumGray] (-6, 3) -- (6, 3);

% رالز (چپ‌تر)
\fill[LeftRedLight, opacity=0.8] (-4, 3) circle (12pt);
\node[text=white, font=\bfseries] at (-4, 3) {\fa{رالز}};
\node[below, text width=4cm, align=center, text=DarkGray, font=\small] at (-4, 2.2) {
    \fa{عدالت = انصاف}\\
    \fa{توزیع عادلانه مهم است}\\
    \fa{دولت رفاه موجه است}
};

% نوزیک (راست‌تر)
\fill[RightBlueLight, opacity=0.8] (4, 3) circle (12pt);
\node[text=white, font=\bfseries] at (4, 3) {\fa{نوزیک}};
\node[below, text width=4cm, align=center, text=DarkGray, font=\small] at (4, 2.2) {
    \fa{عدالت = استحقاق}\\
    \fa{فرآیند کسب مهم است}\\
    \fa{بازتوزیع ناعادلانه است}
};

% فلش مقابله
\draw[thickarrow, color=OrangeAccent] (-2.5, 3) -- (-1, 3);
\draw[thickarrow, color=OrangeAccent] (2.5, 3) -- (1, 3);
\node[text=OrangeAccent, font=\bfseries] at (0, 3) {\fa{⚔}};

% برچسب‌ها
\node[text=LeftRed, font=\small\bfseries] at (-6.5, 3) {\fa{چپ}};
\node[text=RightBlue, font=\small\bfseries] at (6.5, 3) {\fa{راست}};

\end{tikzpicture}
\caption{مقایسه نظریه‌های عدالت رالز و نوزیک}
\label{fig:rawls-nozick}
\end{figure}

\begin{comparebox}[مقایسه: رالز و نوزیک]
\begin{table}[H]
\centering
\begin{tabular}{r|c|c}
\toprule
\textbf{موضوع} & \textbf{رالز} & \textbf{نوزیک} \\
\midrule
تعریف عدالت & الگوی توزیع & فرآیند کسب \\
نابرابری & موجه اگر به نفع فقرا & موجه اگر فرآیند عادلانه \\
مالکیت & حق مشروط & حق مطلق \\
دولت & رفاهی/توزیعی & حداقلی (شب‌پاس) \\
مالیات & ابزار عدالت & شکلی از بردگی \\
جایگاه سیاسی & لیبرالیسم برابری‌خواه & لیبرتارینیسم \\
\bottomrule
\end{tabular}
\end{table}
\end{comparebox}

\subsection{برابری: صوری یا واقعی؟}

یکی از مهم‌ترین اختلافات درباره معنای «برابری» است:

\begin{defbox}[تعریف: انواع برابری]
\begin{itemize}[nosep]
    \item \textbf{برابری صوری (قانونی):} همه در برابر قانون برابرند، حقوق یکسان
    \item \textbf{برابری فرصت:} همه شانس برابر برای رقابت دارند
    \item \textbf{برابری نتیجه:} توزیع نسبتاً برابر درآمد و ثروت
    \item \textbf{برابری پایه:} همه به حداقلی از رفاه دسترسی دارند
\end{itemize}
\end{defbox}

\begin{debatebox}[برابری فرصت کافی است؟]
\textbf{راست (و برخی لیبرال‌ها):}
\begin{itemize}[nosep]
    \item برابری فرصت کافی و عادلانه است
    \item نابرابری نتیجه، پاداش تلاش و استعداد است
    \item برابری اجباری نتیجه، آزادی را نابود می‌کند
\end{itemize}

\textbf{چپ:}
\begin{itemize}[nosep]
    \item برابری فرصت بدون برابری پایه، توهم است
    \item فرزند فقیر و فرزند ثروتمند «فرصت برابر» ندارند
    \item نابرابری شدید، خود فرصت‌ها را نابرابر می‌کند
\end{itemize}
\end{debatebox}

% ══════════════════════════════════════════════════════════════════════════════
\section{آزادی: مفهومی مورد مناقشه}
% ══════════════════════════════════════════════════════════════════════════════

همه از «آزادی» دم می‌زنند، اما منظورشان یکی نیست.

\subsection{آزادی منفی و مثبت: تمایز برلین}

\begin{personbox}[آیزایا برلین (۱۹۰۹-۱۹۹۷)]

\textbf{ملیت:} روسی-بریتانیایی

\textbf{اثر کلیدی:} \book{دو مفهوم آزادی} (۱۹۵۸)

\textbf{ایده محوری:} تمایز آزادی منفی و مثبت

\mediumspace

\person{برلین} دو مفهوم متفاوت از آزادی را شناسایی کرد:

\textbf{آزادی منفی:} آزادی از مداخله دیگران. «چه حوزه‌ای وجود دارد که در آن فرد آنچه می‌خواهد بکند؟»

\textbf{آزادی مثبت:} آزادی برای خودآیینی. «چه کسی یا چه چیزی حاکم بر من است؟»

\mediumspace

\textbf{هشدار برلین:} آزادی مثبت می‌تواند به استبداد بینجامد—وقتی کسی مدعی شود «خود واقعی» شما چیست.
\end{personbox}

\begin{figure}[htbp]
\centering
\begin{tikzpicture}[scale=0.9]

% عنوان
\node[font=\large\bfseries, text=DarkGray] at (0, 6) {\fa{دو مفهوم آزادی}};

% کادر آزادی منفی
\fill[RightBlue, opacity=0.15, rounded corners=10pt] (-7, 0) rectangle (-0.5, 5);
\draw[line width=2pt, color=RightBlue, rounded corners=10pt] (-7, 0) rectangle (-0.5, 5);

\node[text=RightBlue, font=\bfseries\large] at (-3.75, 4.5) {\fa{آزادی منفی}};
\node[text=RightBlue, font=\normalsize] at (-3.75, 3.8) {\fa{Freedom From}};

\node[text=DarkGray, font=\small, text width=5.5cm, align=center] at (-3.75, 2.5) {
    \fa{آزادی از مداخله}\\[5pt]
    \fa{«رهایم بگذارید»}\\[5pt]
    \fa{حوزه خصوصی}\\[5pt]
    \fa{محدودیت قدرت دولت}
};

\node[text=RightBlueDark, font=\small\bfseries] at (-3.75, 0.7) {\fa{لیبرالیسم کلاسیک}};

% کادر آزادی مثبت
\fill[LeftRed, opacity=0.15, rounded corners=10pt] (0.5, 0) rectangle (7, 5);
\draw[line width=2pt, color=LeftRed, rounded corners=10pt] (0.5, 0) rectangle (7, 5);

\node[text=LeftRed, font=\bfseries\large] at (3.75, 4.5) {\fa{آزادی مثبت}};
\node[text=LeftRed, font=\normalsize] at (3.75, 3.8) {\fa{Freedom To}};

\node[text=DarkGray, font=\small, text width=5.5cm, align=center] at (3.75, 2.5) {
    \fa{آزادی برای خودآیینی}\\[5pt]
    \fa{«خودم را محقق کنم»}\\[5pt]
    \fa{خودفرمانی جمعی}\\[5pt]
    \fa{توانمندسازی}
};

\node[text=LeftRedDark, font=\small\bfseries] at (3.75, 0.7) {\fa{جمهوری‌خواهی / سوسیالیسم}};

\end{tikzpicture}
\caption{تمایز آزادی منفی و مثبت (آیزایا برلین)}
\label{fig:two-freedoms}
\end{figure}

\begin{table}[htbp]
\centering
\caption{نمونه‌هایی از کاربرد دو مفهوم آزادی}
\label{tab:freedom-examples}
\begin{tabular}{r|c|c}
\toprule
\textbf{مسئله} & \textbf{آزادی منفی} & \textbf{آزادی مثبت} \\
\midrule
فقر & دولت مداخله نکند & دولت توانمند کند \\
آموزش & انتخاب والدین & دسترسی همگانی \\
بهداشت & بازار آزاد & حق سلامت \\
اشتغال & قرارداد آزاد & حق کار \\
بیان & سانسور نباشد & دسترسی به رسانه \\
\bottomrule
\end{tabular}
\end{table}

\subsection{آزادی جمهوری‌خواهانه: راه سوم؟}

برخی فیلسوفان معاصر (مانند \person{فیلیپ پتیت}) مفهوم سومی از آزادی پیشنهاد کرده‌اند: \keyword{آزادی به مثابه عدم سلطه}.

\begin{defbox}[تعریف: آزادی جمهوری‌خواهانه]
آزادی = نبود سلطه (domination)

انسان آزاد کسی نیست که صرفاً مداخله نمی‌بیند، بلکه کسی است که تحت سلطه کسی نیست—کسی که می‌تواند خودسرانه در زندگی‌اش مداخله کند.

\textbf{مثال:} برده‌ای که ارباب مهربان دارد، مداخله نمی‌بیند اما آزاد نیست—زیرا تحت سلطه است.
\end{defbox}

% ══════════════════════════════════════════════════════════════════════════════
\section{اختلافات درون‌مکتبی}
% ══════════════════════════════════════════════════════════════════════════════

مهم است بدانیم که این اختلافات فلسفی نه‌تنها میان چپ و راست، بلکه \textbf{درون} هر جناح نیز وجود دارد.

\subsection{اختلاف درون چپ}

\begin{debatebox}[چپ: آزادی فردی یا رهایی جمعی؟]
\textbf{چپ آزادی‌خواه (آنارشیست‌ها، برخی سوسیالیست‌ها):}
\begin{itemize}[nosep]
    \item آزادی فردی مقدس است
    \item دولت، حتی دولت کارگری، تهدید آزادی است
    \item رهایی باید از پایین بیاید
\end{itemize}

\textbf{چپ دولت‌گرا (لنینیست‌ها، سوسیال‌دموکرات‌ها):}
\begin{itemize}[nosep]
    \item رهایی جمعی مقدم است
    \item دولت ابزار لازم برای تحقق برابری است
    \item آزادی فردی بدون برابری، بی‌معناست
\end{itemize}
\end{debatebox}

\subsection{اختلاف درون راست}

\begin{debatebox}[راست: سنت یا آزادی بازار؟]
\textbf{محافظه‌کاران سنتی:}
\begin{itemize}[nosep]
    \item سنت، مذهب، خانواده مهم‌تر از بازارند
    \item سرمایه‌داری افسارگسیخته جامعه را ویران می‌کند
    \item دولت باید از ارزش‌ها حمایت کند
\end{itemize}

\textbf{لیبرتارین‌ها / لیبرال‌های اقتصادی:}
\begin{itemize}[nosep]
    \item آزادی فردی و بازار آزاد مقدم‌اند
    \item سنت نباید با زور تحمیل شود
    \item دولت حداقلی بهترین است
\end{itemize}
\end{debatebox}

\begin{figure}[htbp]
\centering
\begin{tikzpicture}[scale=0.8]

% عنوان
\node[font=\large\bfseries, text=DarkGray] at (0, 7) {\fa{تنوع درونی چپ و راست در فلسفه سیاسی}};

% ─────────── چپ ───────────
\fill[LeftRed, opacity=0.15, rounded corners=8pt] (-7.5, 0) rectangle (-0.5, 6);
\draw[line width=2pt, color=LeftRed, rounded corners=8pt] (-7.5, 0) rectangle (-0.5, 6);
\node[text=LeftRed, font=\bfseries\large] at (-4, 5.5) {\fa{چپ}};

% زیرمجموعه‌های چپ
\fill[LeftRedLight, opacity=0.5, rounded corners=5pt] (-7.2, 2.8) rectangle (-4.2, 5);
\node[text=LeftRedDark, font=\small\bfseries] at (-5.7, 4.5) {\fa{آزادی‌خواه}};
\node[text=DarkGray, font=\scriptsize, text width=2.5cm, align=center] at (-5.7, 3.5) {
    \fa{آنارشیست}\\
    \fa{چپ نو}
};

\fill[LeftRedDark, opacity=0.4, rounded corners=5pt] (-4, 2.8) rectangle (-0.8, 5);
\node[text=white, font=\small\bfseries] at (-2.4, 4.5) {\fa{دولت‌گرا}};
\node[text=VeryLightGray, font=\scriptsize, text width=2.5cm, align=center] at (-2.4, 3.5) {
    \fa{لنینیست}\\
    \fa{سوسیال‌دموکرات}
};

% نقطه اختلاف چپ
\node[text=OrangeAccent, font=\small] at (-4, 1.5) {\fa{اختلاف: نقش دولت}};

% ─────────── راست ───────────
\fill[RightBlue, opacity=0.15, rounded corners=8pt] (0.5, 0) rectangle (7.5, 6);
\draw[line width=2pt, color=RightBlue, rounded corners=8pt] (0.5, 0) rectangle (7.5, 6);
\node[text=RightBlue, font=\bfseries\large] at (4, 5.5) {\fa{راست}};

% زیرمجموعه‌های راست
\fill[RightBlueLight, opacity=0.5, rounded corners=5pt] (0.8, 2.8) rectangle (3.8, 5);
\node[text=RightBlueDark, font=\small\bfseries] at (2.3, 4.5) {\fa{محافظه‌کار}};
\node[text=DarkGray, font=\scriptsize, text width=2.5cm, align=center] at (2.3, 3.5) {
    \fa{سنت‌گرا}\\
    \fa{برکی}
};

\fill[CenterGreen, opacity=0.5, rounded corners=5pt] (4, 2.8) rectangle (7.2, 5);
\node[text=white, font=\small\bfseries] at (5.6, 4.5) {\fa{لیبرتارین}};
\node[text=VeryLightGray, font=\scriptsize, text width=2.5cm, align=center] at (5.6, 3.5) {
    \fa{هایک}\\
    \fa{نوزیک}
};

% نقطه اختلاف راست
\node[text=OrangeAccent, font=\small] at (4, 1.5) {\fa{اختلاف: سنت یا بازار؟}};

\end{tikzpicture}
\caption{تنوع فلسفی درون چپ و راست}
\label{fig:internal-diversity-philosophy}
\end{figure}

% ══════════════════════════════════════════════════════════════════════════════
\section{جدول جامع فیلسوفان و مواضعشان}
% ══════════════════════════════════════════════════════════════════════════════

\begin{landscape}
\begin{table}[htbp]
\centering
\caption{فیلسوفان سیاسی کلیدی و مواضعشان}
\label{tab:philosophers-summary}
\small
\begin{tabular}{r|c|c|c|c|c|c}
\toprule
\textbf{فیلسوف} & \textbf{دوره} & \textbf{طبیعت انسان} & \textbf{دولت ایده‌آل} & \textbf{مالکیت} & \textbf{آزادی} & \textbf{جایگاه در طیف} \\
\midrule
هابز & ۱۷ & بدبین & قدرتمند/مطلق & محافظت‌شده & ثانویه & راست اقتدارگرا \\
\midrule
لاک & ۱۷ & معتدل & محدود & حق طبیعی & منفی & لیبرال کلاسیک \\
\midrule
روسو & ۱۸ & خوش‌بین & اراده عمومی & مشکل‌ساز & مثبت & چپ رادیکال \\
\midrule
برک & ۱۸ & محتاط & مشروطه سنتی & سنتی & در چارچوب & محافظه‌کار \\
\midrule
میل & ۱۹ & قابل تربیت & حداقلی لیبرال & خصوصی & منفی (+رشد) & لیبرال \\
\midrule
مارکس & ۱۹ & اجتماعی & زوال دولت & اشتراکی & رهایی جمعی & چپ رادیکال \\
\midrule
رالز & ۲۰ & نامعلوم & رفاهی/توزیعی & مشروط & هر دو & لیبرال چپ \\
\midrule
نوزیک & ۲۰ & فرد محور & حداقلی & مطلق & منفی & لیبرتارین \\
\midrule
برلین & ۲۰ & پیچیده & کثرت‌گرا & خصوصی & منفی (ترجیحاً) & لیبرال \\
\midrule
اوکشات & ۲۰ & نامعلوم & غیرایدئولوژیک & سنتی & منفی & محافظه‌کار \\
\bottomrule
\end{tabular}
\end{table}
\end{landscape}

% ══════════════════════════════════════════════════════════════════════════════
%                         خلاصه و پرسش‌ها
% ══════════════════════════════════════════════════════════════════════════════

\newpage

\begin{summarybox}

\textbf{۱. انسان‌شناسی:} محافظه‌کاری بر نقصان انسان تأکید دارد (هابز)؛ رادیکالیسم بر قابلیت تکامل (روسو). لاک حد میانه‌ای ارائه داد.

\textbf{۲. معرفت‌شناسی:} راست به «خرد جمعی سنت» اعتماد دارد؛ چپ به «عقل فردی و علم». مارکس «نقد ایدئولوژی» را افزود.

\textbf{۳. فلسفه تاریخ:} چپ معمولاً پیشرفت‌گراست؛ محافظه‌کاری محتاط‌تر است. مارکس نظریه «پلکانی» ارائه داد.

\textbf{۴. عدالت:} رالز بر «انصاف» و توزیع تأکید کرد؛ نوزیک بر «استحقاق» و فرآیند. این اختلاف، بنیادی‌ترین شکاف است.

\textbf{۵. آزادی:} برلین دو مفهوم (منفی/مثبت) تشخیص داد. راست معمولاً بر آزادی منفی، چپ بر آزادی مثبت تأکید می‌کند.

\textbf{۶. برابری:} راست به «برابری صوری» قناعت می‌کند؛ چپ «برابری واقعی» می‌خواهد.

\textbf{۷. تنوع درونی:} اختلافات فلسفی درون هر جناح نیز وجود دارد (دولت‌گرا/آزادی‌خواه در چپ؛ سنت‌گرا/لیبرتارین در راست).

\end{summarybox}

\vspace{20pt}

\begin{questionbox}

\begin{enumerate}[nosep, rightmargin=1cm]
    \item آیا می‌توان در انسان‌شناسی خوش‌بین و در سیاست محتاط بود؟ یعنی باور داشت انسان ذاتاً خوب است، اما همچنان از تغییرات رادیکال ترسید؟

    \item رالز و نوزیک هر دو لیبرال‌اند، اما نتیجه‌گیری‌هایشان بسیار متفاوت است. این نشان می‌دهد که درون لیبرالیسم چه طیف وسیعی وجود دارد؟

    \item آیا تمایز آزادی منفی و مثبت واقعاً قاطع است؟ یا هر آزادی، هم جنبه «از» و هم جنبه «برای» دارد؟

    \item مارکس می‌گفت ایده‌ها بازتاب منافع طبقاتی‌اند. آیا این نقد شامل خود نظریه مارکس هم می‌شود؟
\end{enumerate}

\end{questionbox}

% ══════════════════════════════════════════════════════════════════════════════════════